\documentclass[11pt]{article}

\usepackage[margin=1.1in]{geometry}
\usepackage{amsmath,amssymb,amsthm}
\usepackage{stmaryrd}
\usepackage[expansion=false]{microtype}
\usepackage{enumitem}
\usepackage{parskip}
\usepackage[colorlinks=true,linkcolor=black,urlcolor=blue,citecolor=blue]{hyperref}

\newcommand{\Vect}{\mathsf{Vect}}
\providecommand{\Vec}{}\renewcommand{\Vec}{\mathsf{Vec}}
\newcommand{\Set}{\mathsf{Set}}
\newcommand{\Fam}{\mathsf{Fam}}
\newcommand{\Cont}{\mathsf{Cont}}
\newcommand{\Mat}{\mathsf{Mat}}
\newcommand{\Para}{\mathsf{Para}}
\newcommand{\op}{\mathrm{op}}
\newcommand{\fld}{k}
\newcommand{\Hom}{\mathrm{Hom}}
\newcommand{\id}{\mathrm{id}}
\newcommand{\Id}{\mathrm{Id}}
\newcommand{\ext}[1]{\llbracket #1 \rrbracket}
\newcommand{\AttP}{\mathsf{AttP}}
\newcommand{\out}{\mathrm{out}}

\theoremstyle{plain}
\newtheorem{theorem}{Theorem}
\newtheorem{proposition}[theorem]{Proposition}
\newtheorem{corollary}[theorem]{Corollary}
\theoremstyle{definition}
\newtheorem{definition}[theorem]{Definition}

\title{The attention layer is not where the depth is:\\
a stacking law for categorical accounts of attention}
\author{MacBeth \\ \small for Neil}
\date{2026-08-25}

\begin{document}
\maketitle

\begin{abstract}\noindent
Several independent research programmes now recast a single attention (or neural-network)
layer as a piece of structured linear algebra: a parametric span, an endofunctor on
$\Vect$, a discretised Kan extension. They agree on the single layer, and none of them
says anything a transformer engineer does not already know about it. This note argues that
the single layer is the wrong place to look. The mathematics that a practitioner does
\emph{not} already have is the law for how layers \emph{stack}: depth $L$ is a free
hyperparameter in every one of these accounts, and none supplies a compositional law for
it. We give that law. Stacking a linear self-attention layer is the free monoid on its
parameter space --- the tensor algebra $\bigoplus_{L\ge 0}\AttP^{\otimes L}$ --- and this
identification has two consequences a practitioner can use: the residual (skip) connection
is exactly the algebraic \emph{unit} of that free monoid, the degree-$0$ term without which
the construction degenerates; and a trained (data-dependent) depth-$L$ stack, with identity
feature map, is a homogeneous polynomial of degree \emph{exactly} $3^{L}$ in its inputs, a hard, quantitative
obstruction to collapsing a deep attention network into a single weight matrix. The same
degree bound pins the boundary --- linear vs.\ softmax, frozen vs.\ trained --- that the
other accounts each hit qualitatively and leave as a caveat. I flag throughout what is
\emph{proved}, what is \emph{reported} at second hand, and what is \emph{proposal}.
\end{abstract}

\section{The question, and the one that matters}

Category theory has discovered attention several times over. In the last two years at least
three published accounts, developed independently and citing none of the others, take the
same structural step: replace the weight matrix of a layer with a structured, linear-algebra-valued
functor. Vertechi presents a neural-network layer as a \emph{parametric span}
$s\leftarrow P\rightarrow t$ computing one indexed sum
$y[t(p)]\mathrel{+}= x[s(p)]\cdot w[\pi(p)]$~\cite{vertechi}. O'Neill presents a linear
self-attention layer as a \emph{parametric endofunctor} on $\Vect$, a $1$-morphism in
$\Para(\Vect)$~\cite{oneill}. Mahadevan presents a layer as a weighted, discretised
\emph{left Kan extension} --- a coend $\int^{\sigma} W(t,\sigma)\otimes X(\sigma)$ over a
neighbourhood category~\cite{ket}. Three lineages, one move.\footnote{The pattern is wider
than the three accounts I rely on here. At least two further categorical treatments of
attention/deep-learning layers have been reported to me at second hand --- a
presheaf/equivariant account and an applicative-functor account. I have not read them at
source this cycle, so I neither cite nor rely on them; they would only strengthen the point
below, since none of them supplies a stacking law either.}

Each account is correct, and each says the same thing about a single layer: it is a weighted
aggregation, hence a span / a functor / a coend. Here is the uncomfortable fact for anyone
who wants category theory to earn its keep in machine learning: \emph{a transformer engineer
already knows this}. That one layer mixes and re-weights token representations is not news;
naming the mix a coend renames it without adding a usable fact. If the categorical content of
attention were exhausted by the single layer, the honest verdict would be ``true, and not
interesting.''

So we should ask what these accounts \emph{leave out}, in unison. The answer is sharp. Every
one of them models a \emph{single} layer (or a single weighted aggregation) and then treats
the depth $L$ --- the number of stacked layers --- as an experimental hyperparameter, external
to the mathematics. Not one of them gives a compositional law for depth: a statement of what
object a depth-$L$ stack \emph{is}, as a function of $L$, and how the depths compose. That
missing law is the whole subject of this note, because it is the one place where the categorical
picture tells a practitioner something they did not already have.

\paragraph{The result.} Read one linear self-attention layer as the extension of a one-shape
\emph{linear container} $C=(\ast,\AttP)$, where $\AttP=\mathsf{QP}\oplus\mathsf{KP}\oplus\mathsf{VP}$
is O'Neill's parameter space (the weight spaces of $W_Q,W_K,W_V$) and the induced endofunctor is
$F(W)=\AttP\otimes W$. Then the following is a theorem (Section~\ref{sec:stacking}), proved and
machine-checked in the linear/frozen regime.

\begin{theorem}[Stacking law, informal]\label{thm:main-informal}
Stacking linear self-attention layers is the free monoid on $\AttP$ in $(\Vect,\otimes)$: the
depth-graded object
\[
  T(\AttP)\;=\;\bigoplus_{L\ge 0}\AttP^{\otimes L},
\]
the tensor algebra, with layer-composition given by concatenation of parameter tensors. Two
facts about this object are usable:
\begin{enumerate}[label=(\roman*),leftmargin=2.2em]
  \item the residual/skip connection is precisely the \emph{unit} of this free monoid --- the
    degree-$0$ summand $\AttP^{\otimes 0}=\fld$ --- and without it the construction is not the
    free monoid but a single non-terminating tower;
  \item a trained (data-dependent) depth-$L$ stack, with identity feature map, is a homogeneous
    polynomial map of degree \emph{exactly} $3^{L}$ in its input tokens, so it is a single weight
    matrix only at $L=0$.
\end{enumerate}
\end{theorem}

Part~(ii) is the quantitative form of a boundary the other accounts each hit and record as a
caveat: Mahadevan's coend is a genuine enriched Kan extension \emph{only} in the representable
case, and softmax weights are not representable~\cite[\S10.6]{ket}; Sargsyan proves softmax
attention cannot be functorial for non-trivial groups~\cite{sargsyan}. The degree $3^{L}$ makes
``linear collapses, trained does not'' an exact statement rather than a warning.

\paragraph{What is new, and what is not.} The machinery --- that a monoid in a monoidal category
of container-like objects is an enriched category, and that layer-composition is monoidal
composition --- is not mine; it is developed by Dorta, Jarvis and Niu over a general base
category~\cite[Thm.~4.3]{djn},\footnote{\emph{Scope note.} DJN's composition product is
\emph{weighted}: at a position $(i,j\colon A_i\to J)$ its direction object is
$\prod_a\prod_b(u_{i,a}\cdot v_{ja,b})$ (their Def.~3.5 / Lemma~3.6), where ours carries no outer $u$
factor. The two coincide exactly at $\mathcal C=1$; their Day tensor $\otimes$, by contrast, agrees
with ours over any base. Whether the $\Fam(\Vect_{\mathrm{fd}}^{\mathrm{op}})$ comonoid
classification used here is a case of their Theorem~4.3 is therefore an open question, not a settled
subsumption --- the two classifications visibly differ (families of $k$-algebras versus all enriched
categories). We claim nothing either way here.} and the free-monad-as-tensor-algebra fact is classical. The
contribution here is threefold and specific: (a) the \emph{placement} of a published transformer
construction (O'Neill's) inside this machinery, functorially, at the level of unit and
multiplication and not merely objects; (b) a \emph{correction} to that construction, internal to
its own composition rule --- the free monad is the coproduct $\bigoplus_L\AttP^{\otimes L}$, not the
colimit of the layer-powers that O'Neill names --- together with the reading of the residual
connection as the pointing that repairs it; and (c) the degree-$3^L$ boundary, which is a fact
about the trained network, not about the notation. None of this is a restatement of the single
layer; all of it is about how the layers stack.

\paragraph{Outline.} Section~\ref{sec:layer} fixes the concrete linear-attention layer and shows
the three accounts on it, so that the shared gap and the shared wall are visible on one example.
Section~\ref{sec:stacking} proves the stacking law and reads off the residual-as-unit correction.
Section~\ref{sec:degree} proves the degree-$3^L$ boundary and locates the regime --- in-context /
KV-cache --- where depth genuinely does collapse to matrix product. Section~\ref{sec:scorecard}
is an honest scorecard: proved, reported, proposal, and a direct answer to whether any of this
clears the bar of telling a practitioner something new.

\section{One layer, three ways, one wall}\label{sec:layer}

Fix the object of study concretely, so that every abstraction lands on it. A
\emph{linear-attention head} is $H=(\varphi,W_Q,W_K,W_V)$ with a feature map
$\varphi:\mathbb{R}^{d_k}\to\mathbb{R}^{d_\varphi}$ and learned linear maps $W_Q,W_K,W_V$. On a
token sequence $x_1,\dots,x_n$ it forms keys $k_j=W_Kx_j$, values $v_j=W_Vx_j$, and the
\emph{KV-state}
\[
  S \;=\; \sum_{j=1}^{n} \varphi(k_j)\otimes v_j \;\in\; \mathsf{Mat}(\Vec)(d_\varphi,d_v),
\]
a single matrix over $\Vec$. It reads out a query $q$ as $\out(q)=\varphi(W_Qq)^{\!\top}S$. This
is the standard linear (kernelised) attention head; softmax attention replaces the linear
readout by a normalised one, and we return to it in Section~\ref{sec:degree}.

The three categorical accounts each describe this same head.

\begin{description}[leftmargin=1.3em,style=nextline]
\item[Span (Vertechi).] The index pattern $j\mapsto(\text{source }s(j),\text{target }t(j))$ with the
  weight action on positions is a span $s\leftarrow P\rightarrow t$; the head's aggregation
  $y[t(p)]\mathrel{+}=x[s(p)]\cdot w[\pi(p)]$ is its indexed sum~\cite{vertechi}. In container language
  this is a shape/position pair with a linear action on positions --- though Vertechi does not use
  container or polynomial-functor language anywhere; the slides contain the formula and none of the
  vocabulary, which is precisely why the container reading is available to contribute rather than
  already claimed.
\item[Endofunctor (O'Neill).] Bundling the weight spaces as $\AttP=\mathsf{QP}\oplus\mathsf{KP}\oplus\mathsf{VP}$,
  the head is a parametric $1$-morphism $(\AttP,\mathrm{att}):X\to Y$ in $\Para(\Vect)$ with
  $\mathrm{att}:\AttP\otimes X\to Y$ linear in both arguments, inducing an endofunctor
  $F(W)=\AttP\otimes W$ on $\Vect$~\cite[Thm.~3.1]{oneill}.
\item[Kan extension (Mahadevan).] Taking the query's neighbourhood as an index category, the readout
  is a weighted left Kan extension --- a coend $\int^{\sigma}W(t,\sigma)\otimes X(\sigma)$ --- with
  ordinary attention the singleton-neighbourhood case~\cite{ket}.
\end{description}

Two observations, and they are the point of the section.

\emph{None of the three gives a stacking law.} Each describes one layer. O'Neill alone even asks how
layers compose, and answers with a free monad on $F$ --- correctly in intent, but with the wrong
diagram named, as Section~\ref{sec:stacking} shows. Vertechi and Mahadevan leave depth entirely
outside the formalism. So ``how does a depth-$L$ stack depend on $L$?'' is unanswered in all three.

\emph{All three hit the same wall.} The clean categorical identity holds only in a linear /
representable / frozen regime, and each account records the wall as a caveat rather than measuring it.
Mahadevan is explicit: the coend is a genuine enriched Kan extension only when the weights are
representable, and softmax is not representable, so generic attention is only ``Kan-style,'' an
analogy and not an identity~\cite[\S10.6]{ket}. O'Neill declares the layer linear and defers softmax.
Sargsyan, from an entirely different direction, proves softmax attention cannot be functorial for
non-trivial groups~\cite{sargsyan}. The field keeps rediscovering one boundary, qualitatively.

Neither observation, on its own, clears the bar of telling a practitioner something new: that a single
attention layer is a coend, or is linear-only-until-softmax, is known. They set up the gap. The next two
sections fill it, and that is where the new content is.

\section{The stacking law: depth is a free monoid, and the residual is its unit}\label{sec:stacking}

\paragraph{The container reading.} A \emph{linear container} is a pair $(S,(P_s)_{s\in S})$ --- a set
$S$ of shapes, a finite-dimensional space $P_s$ at each --- with extension the endofunctor
$\ext{S,P}\,W=\bigoplus_s \Hom(P_s,W)$ on $\Vect$~\cite{vcont}. One shape over the space $A$ gives
$\ext{\ast,A}\,W=\Hom(A,W)$; in finite dimension this is the functor $A^{*}\!\otimes W$, and, taking the
covariant reading faithful to O'Neill's $\mathrm{att}$, we write $F(W)=A\otimes W$. Containers compose by
substitution, $\ext{C\lhd D}=\ext{C}\circ\ext{D}$, and the free monoid for $\lhd$ --- the free
$\lhd$-monoid on a container --- is the container-theoretic free monad~\cite{djn}. Set $C=(\ast,\AttP)$,
so $F(W)=\AttP\otimes W$ is exactly O'Neill's layer endofunctor.

The one fact that drives everything is O'Neill's own composition rule. In $\Para(\Vect)$, parametric
morphisms $(P,f)$ and $(Q,g)$ compose to $(Q\otimes P,\,h)$: \emph{parameters tensor}. Hence a depth-$L$
stack has parameter space $\AttP^{\otimes L}$~\cite[\S3]{oneill}. This is precisely container composition.

\begin{proposition}[Composition is $\lhd$]\label{prop:comp}
For finite-dimensional $A,B$, $(\ast,A)\lhd(\ast,B)=(\ast,A\otimes B)$, and the one-shape linear containers
under $\lhd$ are monoidally isomorphic to $(\Vect_{\mathrm{fd}},\otimes,\fld)$, with unit container
$(\ast,\fld)=\Id$.
\end{proposition}

\begin{proof}
$\ext{\ast,A}\circ\ext{\ast,B}\,(W)=\Hom(A,\Hom(B,W))\cong\Hom(A\otimes B,W)=\ext{\ast,A\otimes B}\,W$ by the
tensor--hom adjunction (currying), naturally in $W$. The unit: $\ext{\ast,\fld}\,W=\Hom(\fld,W)=W$, and
$A\otimes\fld\cong A\cong\fld\otimes A$. Associativity and unit coherence transport from $(\Vect,\otimes)$.
\end{proof}

So O'Neill's parameter-tensor composition \emph{is} the container $\lhd$, and a depth-$L$ stack is the
$L$-fold power $(\ast,\AttP)^{\lhd L}=(\ast,\AttP^{\otimes L})$. Now take the free monoid.

\begin{theorem}[The stacking object]\label{thm:stack}
The free $\lhd$-monoid on $C=(\ast,\AttP)$ is the container $C^{*}=(\mathbb{N},(\AttP^{\otimes L})_{L\in\mathbb{N}})$:
its shapes are layer-counts, the position space of shape $L$ is $\AttP^{\otimes L}$. Its extension is
\[
  \ext{C^{*}}\,W \;=\; \bigoplus_{L\ge 0}\AttP^{\otimes L}\otimes W \;=\; T(\AttP)\otimes W,
  \qquad T(\AttP)=\bigoplus_{L\ge 0}\AttP^{\otimes L},
\]
the tensor algebra, with unit $\eta$ the degree-$0$ inclusion $W\cong\AttP^{\otimes 0}\otimes W$ and
multiplication $\mu$ induced by concatenation $\AttP^{\otimes m}\otimes\AttP^{\otimes n}\to\AttP^{\otimes(m+n)}$.
This $(T,\eta,\mu)$ is O'Neill's free monad on $F$, at the level of $(\eta,\mu)$ and not merely objects.
\end{theorem}

\begin{proof}
By Proposition~\ref{prop:comp} the one-shape linear containers under $\lhd$ are $(\Vect_{\mathrm{fd}},\otimes)$
with $(\ast,\AttP)$ at the object $\AttP$; the free $\lhd$-monoid on it is the free monoid on $\AttP$ in
$(\Vect,\otimes)$, which is the tensor algebra $T(\AttP)$ with concatenation and degree-$0$ unit
(classical: the tensor algebra is the free associative algebra). As a container this is
$(\mathbb{N},(\AttP^{\otimes L}))$ since the $L$-fold $\lhd$-power is $(\ast,\AttP^{\otimes L})$
(Proposition~\ref{prop:comp}, induction) and the free monoid coproduces the powers. A monoid $M$ in a
monoidal category yields the monad $M\otimes(-)$ with the monoid laws as monad laws; the universal property
of the free monad on $M\otimes(-)$ coincides with that of the free monoid on the generating object, so the
structures agree, not just the underlying functors. The general free-$\lhd$-monoid multiplication is
tree-grafting~\cite{djn}; here $C$ has one shape, so its trees are paths and grafting degenerates to list
concatenation, which is $\mu$.
\end{proof}

The engineer's payoff is in reading the unit correctly, and it also corrects the literature.

\begin{proposition}[Residual is the unit; O'Neill names the wrong diagram]\label{prop:residual}
Let $F(W)=\AttP\otimes W$ with $\dim\AttP=p\ge 2$. The free monad $\mathrm{Free}(F)$ is the coproduct
$\bigoplus_{L\ge 0}F^{L}$ (equivalently the colimit of the Adámek chain of partial sums
$Y_{n}=\bigoplus_{L<n}\AttP^{\otimes L}\otimes W$), \emph{not} the colimit of the chain of bare powers
$\id\to F\to F^{2}\to\cdots$ that~\cite[Thm.~3.2]{oneill} names. The two diagrams have different colimits:
the partial-sum stage has dimension $w\,(p^{n}-1)/(p-1)$ (a geometric \emph{series}), the bare-power stage
$w\,p^{n}$ (a single growing \emph{term}); for $p=2,\,w=3$ these are $[3,9,21,45,\dots]$ versus
$[3,6,12,24,\dots]$, equal only at the base. Moreover $F$ is pointable but never well-pointed
($a\otimes u\otimes w\ne u\otimes a\otimes w$ for $a\ne 0$, $p\ge 2$), so the colimit-of-powers theorem does
not apply. The pointing that repairs the chain is exactly the residual connection: writing a real block as
$x\mapsto x+\mathrm{att}(\theta,x)$ gives the pointed endofunctor $F'(W)=W\oplus\AttP\otimes W$, whose unit
$\eta=\iota_{1}:W\to F'(W)$ is the degree-$0$ ``stop'' summand, and $\mathrm{Free}(F')=\mathrm{Free}(F)=T(\AttP)\otimes(-)$.
\end{proposition}

\begin{proof}
$\AttP\otimes(-)$ is finitary on $\Vect$ (a left adjoint, preserving filtered colimits), so by Adámek's
theorem the initial algebra of $Y\mapsto W\oplus F(Y)$ --- that is, $\mathrm{Free}(F)(W)$ --- is the colimit
of $Y_{0}=0$, $Y_{n+1}=W\oplus F(Y_{n})$. Computing, $Y_{1}=W$, $Y_{2}=W\oplus\AttP\otimes W$, and inductively
$Y_{n}=\bigoplus_{L<n}\AttP^{\otimes L}\otimes W$ with the partial-sum inclusions as transitions; the colimit
is $\bigoplus_{L\ge 0}\AttP^{\otimes L}\otimes W$, matching Theorem~\ref{thm:stack}. The bare-power chain has
$n$-th object the single power $\AttP^{\otimes n}\otimes W$; its stage-dimensions $w\,p^{n}$ diverge from the
partial sums for every $n\ge 1$, so the diagrams and their colimits differ. Pointability and the failure of
well-pointedness are the displayed identities. The pointed $F'$ is the two-shape container
$(\{\text{stop},\text{layer}\},\,P_{\text{stop}}=0,\,P_{\text{layer}}=\AttP)$, whose empty-position stop shape
is the degree-$0$ summand turning the free monoid into the sum over layer-counts rather than one tower.
\end{proof}

There is a usable sentence here for someone who builds these networks: the residual connection is not
decoration and not only an optimisation trick; it is the algebraic unit that makes a depth-$L$ stack a
\emph{sum over depths} $0,1,\dots,L$ rather than a single depth-$L$ tower. Remove it and the free-monoid
structure --- the thing that lets you reason about stacks of every depth at once --- is gone. O'Neill's
intended object is right (his own $Q\otimes P$ rule forces the parameter space $\AttP^{\otimes L}$); the
one-line fix is to take the coproduct over $L$, which is by construction the free $\lhd$-monoid. These
three $\lhd$-monoid laws are formalised and machine-checked, sorry-free, in Lean~\cite{leanfree}.

\section{The degree boundary: why a trained stack is not one matrix}\label{sec:degree}

The stacking law of Section~\ref{sec:stacking} is exact in the linear, frozen-parameter regime O'Neill works
in --- where a layer is the parameter-tensoring functor $\AttP\otimes(-)$. The moment the parameters (the
KV-state) become functions of the input, as in genuine self-attention, a single clean invariant separates the
trained network from any single matrix.

\begin{theorem}[Degree of a live stack]\label{thm:degree}
With $\varphi=\id$, one live self-attention layer sends the token matrix $X\in\mathbb{R}^{n\times d}$ to
outputs $\out_{i}=\bigl(W_{V}X^{\!\top}X\,W_{K}^{\!\top}W_{Q}\bigr)x_{i}$, homogeneous of degree $3$ in the
entries of $X$. A live depth-$L$ stack is homogeneous of degree \emph{exactly} $3^{L}$.
\end{theorem}

\begin{proof}
$S=\sum_{j}k_{j}v_{j}^{\!\top}=W_{K}\bigl(\sum_{j}x_{j}x_{j}^{\!\top}\bigr)W_{V}^{\!\top}=W_{K}X^{\!\top}X\,W_{V}^{\!\top}$
is degree $2$ in $X$; then $\out_{i}=S^{\!\top}(W_{Q}x_{i})=W_{V}X^{\!\top}X\,W_{K}^{\!\top}W_{Q}x_{i}$ multiplies
a degree-$2$ factor by a degree-$1$ factor, giving degree $3$, every monomial with exactly three $X$-factors.
Induction on $L$: if the depth-$(L-1)$ output $Y$ is homogeneous of degree $3^{L-1}$, then layer $L$ forms
$W_{V}Y^{\!\top}Y\,W_{K}^{\!\top}W_{Q}Y$, of degree $2\cdot 3^{L-1}+3^{L-1}=3^{L}$.
\end{proof}

\begin{corollary}[No collapse]\label{cor:nocollapse}
A live depth-$L$ stack with $L\ge 1$ is not any single matrix acting on its input: a matrix acts in degree $1$,
the stack in degree $3^{L}\ge 3$, and degree is an invariant. This is the quantitative reason O'Neill must model
stacking by a non-collapsing free monad rather than one morphism: the depths have genuinely different degrees, so
no single object can carry the whole stack.
\end{corollary}

This is worth stating plainly to a practitioner, because the opposite is a common informal belief. Linear
attention is often summarised as ``just a matrix $S$'' --- and in one precise regime it is. Fix the context and
vary only the query (the in-context / KV-cache regime): the head is literally the matrix $u\mapsto u^{\!\top}S$,
and then two operations are genuinely functorial. Contexts compose by $\oplus$ --- appending tokens adds
KV-states, $S(C\cdot C')=S(C)\oplus S(C')$, which is exactly linear attention in its recurrent (linear-RNN)
form. Depths compose by $\odot$, the $\Mat(\Vec)$ product $(P\odot Q)(a,c)=\bigoplus_{b}P(a,b)\otimes Q(b,c)$:
stacking frozen heads is honest matrix multiplication. So in the KV-cache regime the ``pipeline = matrix product''
picture is a theorem. Corollary~\ref{cor:nocollapse} says exactly where it stops: the moment the KV-state is a
function of the input --- the moment you are looking at a \emph{trained} network rather than a frozen context ---
the degree climbs $1\to 3\to 9\to\cdots\to 3^{L}$ and the collapse to one matrix is impossible.

Softmax lies even further out. The softmax readout $\sum_{j}\mathrm{softmax}_{j}(\langle q,k_{j}\rangle)v_{j}$ is
not $u^{\!\top}S$ for any fixed $S$: it is not homogeneous of degree $1$ (the normaliser depends nonlinearly on
$q$) and it is bounded (a convex combination of the $v_{j}$) where $u^{\!\top}S$ is unbounded. It is not a
$\odot$-composition at all. Sargsyan reaches the same conclusion at a categorical level: softmax attention cannot
be functorial for non-trivial groups~\cite{sargsyan}, and Mahadevan's representability caveat~\cite[\S10.6]{ket} is
the same wall seen through Kan extensions. The degree-$3^{L}$ statement is the arithmetic behind all three:
representable / linear / frozen is degree $1$ and collapses; trained is degree $3^{L}$ and does not; softmax is not
polynomial at all.

\section{Honest scorecard}\label{sec:scorecard}

I separate what I have established from what I have not, in the register Neil asked for.

\paragraph{Proved (and, where noted, machine-checked).}
\begin{itemize}[leftmargin=1.4em,itemsep=1pt]
\item O'Neill's free monad on the linear self-attention endofunctor is the free $\lhd$-monoid on the one-shape
  linear container $(\ast,\AttP)$, i.e.\ the tensor algebra $\bigoplus_{L}\AttP^{\otimes L}$, functorially in
  $(\eta,\mu)$ (Prop.~\ref{prop:comp}, Thm.~\ref{thm:stack}).
\item The correction: the free monad is the coproduct $\bigoplus_{L}F^{L}$ (Adámek partial sums), not the colimit
  of bare layer-powers; $F$ is pointable but not well-pointed; the residual connection is the pointing / unit
  (Prop.~\ref{prop:residual}). The three $\lhd$-monoid laws are formalised sorry-free in Lean~\cite{leanfree}.
\item A live depth-$L$ linear-attention stack has polynomial degree exactly $3^{L}$; hence no collapse to a single
  matrix for $L\ge 1$ (Thm.~\ref{thm:degree}, Cor.~\ref{cor:nocollapse}). The in-context / KV-cache regime is
  exactly where depth does collapse, to the $\Mat(\Vec)$ product $\odot$.
\end{itemize}

\paragraph{Reported, not verified this cycle.} Two further categorical accounts of attention/DL layers exist that
I have only at second hand and therefore do not cite or lean on. Both, as reported, model a single layer and
neither supplies a stacking law, so they fit the pattern; confirming this needs a full-text read.

\paragraph{Proposal, not proved.} The full unification: that \emph{all} of these accounts are instances of linear
containers over $\Fam(\Vec^{\op})$, with layer-stacking uniformly the $\lhd$-monoid free monad. O'Neill is proved
inside (above). Vertechi's parametric spans are, structurally, one-shape linear containers with a linear action on
positions, and the identification is at the computed level, not re-derived here. Mahadevan's coend is the
representable case of the same construction --- his \S10.6 caveat is, on this reading, the statement that the
identity holds exactly on the linear/representable rung, matching the degree-$1$ vs.\ degree-$3^{L}$ split --- but I
state this as a thesis with evidence, not a theorem. The value of the note does not depend on the full unification
landing: the depth-composition law and the degree bound are the uncontested delta regardless of how many single-layer
lineages ultimately fit under one roof.

\paragraph{Does it clear the bar?} The bar is whether any of this tells someone who builds transformers a fact they
did not have. The single-layer identifications do not, and I have not pretended otherwise --- ``attention is a coend''
is true and known. Two statements here are, I believe, on the right side of the line. First: the residual connection
is the algebraic unit of the free monoid that stacking forms, the degree-$0$ term without which a stack is not a sum
over depths --- a structural role for skip connections that is not the usual optimisation story. Second: a trained
depth-$L$ linear-attention stack is a degree-$3^{L}$ homogeneous map, an exact obstruction to the ``deep attention is
one big matrix'' intuition, with the frozen KV-cache regime pinned as the precise exception. If those two do not
clear the bar, the note should go in the drawer; I do not think they are re-labellings, because each is a
quantitative or structural fact about the trained network, not about the choice of category.

\begin{thebibliography}{9}

\bibitem{oneill}
G.~O'Neill,
\emph{Self-Attention as a Parametric Endofunctor: A Categorical Framework},
arXiv:2501.02931, 2025. (Thm.~3.1: a linear head is a $\Para(\Vect)$ $1$-morphism with parameter space
$\AttP=\mathsf{QP}\oplus\mathsf{KP}\oplus\mathsf{VP}$; \S3: composition tensors parameters; Thm.~3.2: stacking as a
free monad.)

\bibitem{vertechi}
P.~Vertechi,
\emph{Neural network layers as parametric spans} (lecture slides). Core formula
$y[t(p)]\mathrel{+}=x[s(p)]\cdot w[\pi(p)]$; the slides use no container or polynomial-functor vocabulary.

\bibitem{ket}
S.~Mahadevan,
\emph{Kan Extension Transformers: A Categorical Unification of Attention}, arXiv:2605.27259, 2026. (\S10.6: the
coend is a genuine enriched Kan extension only in the representable case; softmax is non-representable.)

\bibitem{sargsyan}
K.~Sargsyan,
\emph{Functorial Neural Architectures from Higher Inductive Types}, arXiv:2603.16123, 2026. (Softmax attention is
not functorial for non-trivial groups.)

\bibitem{djn}
A.~Dorta, M.~Jarvis, N.~Niu,
\emph{Monoidal Structures on Generalized Polynomial Categories}, arXiv:2305.05655. (Thm.~4.3: comonoids for
\emph{their} weighted substitution product over a general base are enriched categories; the
$\lhd$-monoid / free-monad machinery. See the scope note in \S1: their substitution product
coincides with ours only at $\mathcal C=1$; their Day tensor coincides with ours over any base.)

\bibitem{acu}
D.~Ahman, T.~Uustalu,
\emph{Directed Containers as Categories}, EPTCS 207, 2016. (Directed containers $=$ small categories.)

\bibitem{vcont}
MacBeth,
\emph{Linear containers for prompt/response systems} and \emph{Containers over $\Vec$} (Kodamai internal notes,
2026). Definition of linear containers over $\Fam(\Vec^{\op})$ and their extensions.

\bibitem{leanfree}
MacBeth,
\emph{\texttt{Free.lean}} (Kodamai Lean development, 2026). The three free-$\lhd$-monoid laws, sorry-free.

\end{thebibliography}

\end{document}
